\documentclass[journal=jctcce,manuscript=article]{achemso}

\usepackage{amsmath}
\usepackage{amssymb}
\usepackage{booktabs}
\usepackage{graphicx}
\usepackage{siunitx}

\newcommand{\quick}{QUICK}
\newcommand{\keyword}[1]{\texttt{#1}}
\newcommand{\esp}{ESP}
\newcommand{\efield}{EFIELD}
\newcommand{\efg}{EFG}

\SectionNumbersOn

\title{Application Note: On-the-Fly Electrostatic Potentials, Electric Fields, and Field Gradients for Microdroplet Chemistry in QUICK}

\author{Etienne Palos}
\affiliation{Department of Chemistry, Michigan State University, East Lansing, Michigan, United States}
\alsoaffiliation{TODO: Add complete affiliation list}
\email{TODO: email}

\author{TODO: Coauthors}
\affiliation{TODO: Affiliations}

\keywords{electrostatic potential, electric field, electric field gradient, microdroplet chemistry, air-water interface, ab initio molecular dynamics, trajectory analysis, Gaussian basis sets, QUICK}

\begin{document}

\begin{abstract}
Local electrostatic observables provide a compact bridge between electronic
structure and chemically interpretable molecular environments in
microdroplets, sea-spray aerosols, and air-water interfaces.  The
electrostatic potential (ESP), electric field (EFIELD), and electric
field-gradient tensor (EFG) are the scalar, vector, and tensor descriptors
obtained from a single quantum-mechanical charge distribution at arbitrary
probe points.  Here we describe an application-oriented implementation in
\quick{} for evaluating ESP, EFIELD, and analytic EFG on user-supplied
external grids, electron-density isosurfaces, and fixed probe sets along
molecular trajectories.  The target application is trajectory-resolved
microdroplet chemistry, where interfacial charge separation, partial
solvation, photochemical channels, and heterogeneous local fields have
been proposed to control reactivity.  The same post-SCF property engine is
exposed through a no-I/O library API, allowing external molecular dynamics
drivers or trajectory-analysis tools to request ESP/EFIELD/EFG immediately
after each converged single-point, geometry-optimization, or
Born-Oppenheimer molecular dynamics step.  The tensor convention is
explicitly defined as
\(E_i(\mathbf{C})=-\partial V(\mathbf{C})/\partial C_i\) and
\(G_{ij}(\mathbf{C})=\partial E_i(\mathbf{C})/\partial C_j\), the negative
of the electrostatic-potential Hessian convention sometimes also denoted
EFG.  CPU serial and MPI implementations are available for ESP, EFIELD,
and analytic EFG; the existing GPU ESP path defines the accelerator
template for EFIELD and EFG kernels.  File-based and no-I/O API
evaluations agree to better than \(5\times10^{-9}\) a.u. for a water
RHF/STO-3G validation grid, and analytic EFG agrees with a central
finite-difference EFIELD reference to approximately \(10^{-9}\) a.u. in
the current regression case.  We outline an on-the-fly AIMD workflow for
converting electronic trajectories into time series of local potentials,
fields, field magnitudes, tensor traces, Frobenius norms, and
anisotropies at chemically selected probes.
\end{abstract}

\section{Introduction}

Electrostatic observables are among the most direct ways to interpret a
quantum-mechanical density in real space.  The electrostatic potential
reports the scalar potential experienced by a test charge, the electric
field reports the directional force per unit charge, and the electric
field gradient reports how rapidly that field changes locally.  These
quantities are useful for charge-model construction, field-sensitive
reactivity, vibrational Stark interpretation, interfacial electrostatics,
and polarizable embedding.

The motivating application for the present work is microdroplet chemistry.
Aqueous microdroplets and sea-spray aerosols can display chemical behavior
that differs strongly from bulk solution, including interfacial oxidation,
photochemistry, and rate acceleration.  Recent studies, including work by
Francisco and co-workers on sea-spray microdroplet oxidation and hydrogen
peroxide photolysis, have emphasized the possible roles of the air-water
interface, charge separation, partial solvation, and heterogeneous local
electric fields.\cite{RaoFranciscoThiol2023,RaoFranciscoH2O22023}
Recent reviews and perspectives have likewise identified interfacial
charge and local electrostatics as central, but not yet fully resolved,
ingredients in microdroplet reactivity.\cite{LaCourHeadGordon2025}
This creates a practical need for quantum-mechanical tools that can report
not only energies and structures, but also the instantaneous scalar,
vector, and tensor electrostatic environment experienced by reactive
sites.

Trajectory-resolved electrostatics are especially useful because many
chemical environments are intrinsically dynamic.  Recent molecular
simulation studies of interfacial and condensed-phase electrostatics have
shown that analyzing local fields along trajectories can reveal
heterogeneous microscopic environments that are obscured by structural
averages alone.\cite{MartinsCostaRuizLopez2023}  For ab initio molecular
dynamics (AIMD) or Born-Oppenheimer molecular dynamics (BOMD), the natural
computational object is the converged electronic density at each frame.
Once that density exists, evaluating ESP/EFIELD/EFG at selected probe
points is a post-SCF property task, not a new electronic-structure problem.

This Application Note describes the current \quick{} workflow for local
electrostatic analysis along trajectories, with microdroplet and
air-water-interface chemistry as the target use case.  The implementation
is staged: external-grid and electron-density-surface ESP/EFIELD/EFG are
available in the standard input/output path; the same kernels are exposed
through a no-I/O API for trajectory and driver use; and GPU ESP already
exists as the accelerator template for future GPU EFIELD and EFG kernels.
We focus on the application interface, sign conventions, validation
protocol, and data products needed to use these observables in AIMD-style
studies of reactive interfacial environments.

\section{Definitions}

For a probe point \(\mathbf{C}\), nuclear positions \(\mathbf{R}_A\), and
electron density \(\rho(\mathbf{r})\), the electrostatic potential is
\begin{equation}
  V(\mathbf{C}) =
  \sum_A \frac{Z_A}{|\mathbf{C}-\mathbf{R}_A|}
  -
  \int
  \frac{\rho(\mathbf{r})}{|\mathbf{C}-\mathbf{r}|}
  d\mathbf{r}.
  \label{eq:esp}
\end{equation}
The electric field is
\begin{equation}
  E_i(\mathbf{C}) =
  -\frac{\partial V(\mathbf{C})}{\partial C_i},
  \label{eq:efield}
\end{equation}
and the field-gradient tensor reported by \quick{} is
\begin{equation}
  G_{ij}(\mathbf{C}) =
  \frac{\partial E_i(\mathbf{C})}{\partial C_j}
  =
  -\frac{\partial^2 V(\mathbf{C})}{\partial C_i\partial C_j}.
  \label{eq:efg}
\end{equation}
Equation~\ref{eq:efg} is the field-gradient convention.  Some
quantum-chemistry texts define EFG as the Hessian of the electrostatic
potential itself; that convention differs by an overall sign.

In an AO basis, the electronic EFG contribution can be written as
\begin{equation}
  G_{ij}^{\mathrm{elec}}(\mathbf{C}) =
  \sum_{\mu\nu} P_{\mu\nu}
  \frac{\partial^2}{\partial C_i\partial C_j}
  (\mu|r_C^{-1}|\nu),
  \label{eq:ao_efg}
\end{equation}
where \(P_{\mu\nu}\) is the density matrix and
\begin{equation}
  (\mu|r_C^{-1}|\nu) =
  \int
  \chi_\mu(\mathbf{r})
  \frac{1}{|\mathbf{r}-\mathbf{C}|}
  \chi_\nu(\mathbf{r})d\mathbf{r}.
\end{equation}
The analytic CPU EFG implementation evaluates these center derivatives by
Obara--Saika-style recurrence over Gaussian basis functions.\cite{ObaraSaika1986}
A central finite-difference EFG from the analytic EFIELD remains available
as a reference path.

\section{Implementation in QUICK}

\quick{} now supports three complementary property routes:
\begin{enumerate}
  \item File-based external-grid properties through the grid keywords
    for ESP, EFIELD, and EFG.
  \item Electron-density isosurface properties through the corresponding
    surface keywords for ESP, EFIELD, and EFG.
  \item A no-I/O API routine, \keyword{getQuickOEPROP}, that evaluates
    ESP, EFIELD, and/or EFG after a converged API energy or gradient call.
\end{enumerate}

The no-I/O API has the form
\begin{verbatim}
call getQuickOEPROP(npoints, probe_xyz_bohr, ierr, &
                    esp=esp, efield=efield, efg=efg)
\end{verbatim}
where the probe coordinates are in bohr and the optional output arrays are
\begin{verbatim}
esp(npoints)
efield(3,npoints)
efg(3,3,npoints)
\end{verbatim}
The API delegates all property evaluation to the same OEPROP kernels used
by the file-based routes.  It returns a clean error when called before a
converged density exists, when no output property is requested, or when the
probe set is invalid.

The implementation uses the existing shell-pair decomposition in serial
and MPI builds.  For MPI, each rank accumulates its owned electronic
shell-pair contribution and the electronic property arrays are reduced over
\(N_{\mathrm{grid}}\), \(3N_{\mathrm{grid}}\), or \(9N_{\mathrm{grid}}\)
double-precision values for ESP, EFIELD, or EFG, respectively.  The direct
nuclear terms are inexpensive and are evaluated on the host.

\section{Trajectory and AIMD Workflow}

The intended AIMD use pattern is post-SCF and frame local:
\begin{enumerate}
  \item Converge the electronic structure for the current geometry.
  \item Define probe points in a fixed laboratory frame, in a
    molecule-attached local frame, or on a generated density isosurface.
  \item Evaluate ESP/EFIELD/EFG at those points using the file path or
    \keyword{getQuickOEPROP}.
  \item Append compact descriptors to a trajectory CSV or binary analysis
    file.
  \item Continue to the next geometry or dynamics step.
\end{enumerate}
For externally generated trajectories, this workflow is trajectory
post-processing.  For an external MD or QM/MM driver linked against the
\quick{} API, the same call is on-the-fly because properties are requested
immediately after each SCF or gradient evaluation.  A native \quick{} BOMD
driver can use the same hook once the integration, restart, and energy
conservation infrastructure is in place.

The validation trajectory driver writes one row per frame and probe point.
The current CSV columns are:
\begin{verbatim}
frame,probe,energy_hartree,x_ang,y_ang,z_ang,esp,
efield_x,efield_y,efield_z,efield_mag,
efg_xx,efg_xy,efg_xz,efg_yx,efg_yy,efg_yz,
efg_zx,efg_zy,efg_zz,efg_trace,
efg_frobenius,efg_anisotropy
\end{verbatim}
The tensor anisotropy is reported as the rotationally invariant deviatoric
norm,
\begin{equation}
  \eta_G =
  \left[
    \frac{3}{2}
    \sum_{ij}
    \left(G_{ij}-\frac{\mathrm{Tr}(G)}{3}\delta_{ij}\right)^2
  \right]^{1/2}.
  \label{eq:anisotropy}
\end{equation}
This descriptor is useful for detecting locally anisotropic electrostatic
environments, while \(|\mathbf{E}|\), \(\mathrm{Tr}(G)\), and
\(\|G\|_F\) provide complementary scalar summaries.

\section{Validation}

The present validation set covers file-based properties, no-I/O API
properties, trajectory output, analytic EFG, and serial/MPI consistency.
The no-I/O API was compared against the standard file-producing grid
property route for ESP, EFIELD, and EFG.  The test case is water
RHF/STO-3G with four fixed probe points.  The first frame of a three-frame
trajectory driver test is identical to the validation geometry, providing
an additional trajectory-vs-API comparison.

\begin{table}[htbp]
\centering
\caption{Current no-I/O OEPROP API validation for water RHF/STO-3G.}
\begin{tabular}{lcc}
\toprule
Comparison & Values & Max. abs. dev. / a.u. \\
\midrule
File ESP vs API ESP & 4 & \(3.66\times10^{-11}\) \\
File EFIELD vs API EFIELD & 12 & \(4.71\times10^{-9}\) \\
File EFG vs API EFG & 36 & \(3.85\times10^{-9}\) \\
Trajectory frame 1 ESP vs validation & 4 & \(0.0\) \\
Trajectory frame 1 EFIELD vs validation & 12 & \(0.0\) \\
Trajectory frame 1 EFG vs validation & 36 & \(0.0\) \\
\bottomrule
\end{tabular}
\label{tab:api_validation}
\end{table}

Analytic EFG is also checked against the retained numerical EFG reference,
which uses a central finite difference of EFIELD.  In the current external
grid regression case, analytic and numerical EFG agree to approximately
\(10^{-9}\) a.u.  Independent PySCF cross-code validation on external-grid
and density-surface points gives maximum absolute differences of
\(4.2\times10^{-6}\) a.u. or smaller across ESP, EFIELD, and EFG test
sets, with stricter \(10^{-7}\)--\(10^{-8}\) a.u. agreement for the
water RHF/STO-3G surface cases.

\section{GPU and Performance Considerations}

The current implementation separates the property workflow from the SCF
procedure: the density is converged once per frame, and ESP/EFIELD/EFG are
then evaluated on the requested probe set.  For a small number of probe
points, the incremental cost is typically dominated by the SCF.  For large
external grids or dense molecular surfaces, the OEPROP contraction becomes
a natural accelerator target.

\quick{} already contains a CUDA/HIP electronic ESP path for external-grid
OEPROP calculations.  That path uploads grid coordinates and the converged
density, evaluates the electronic attraction contribution on the
accelerator, downloads the electronic property buffer, and adds the direct
nuclear contribution on the host.  The planned GPU EFIELD and EFG kernels
preserve the same host/device division, buffer ordering, and MPI reduction
lengths as the CPU reference.  The acceptance target for the accelerator
extension is CPU/GPU agreement for ESP, EFIELD, analytic EFG, and the
numerical EFG reference on the existing OEPROP regression suite.

\section{Microdroplet-Chemistry Target}

The immediate application target is compact local-electrostatics analysis
of microdroplet, aerosol, and air-water-interface trajectories.  In this
setting, the quantities of interest are not only global dipoles or average
surface potentials.  The chemically relevant questions are local: what
field does a peroxide O--O bond, thiol S--H bond, adsorbed ion, radical
precursor, or surface hydroxyl group experience at the instant a reactive
configuration appears?  How anisotropic is the field around the reactive
center?  Does the local environment look bulk-like, interface-like, or
charge-separated?

The present workflow is designed to answer these questions without
claiming that \quick{} is already a complete droplet dynamics package.
One may generate AIMD or BOMD frames for finite water clusters,
nanodroplets, slabs, or reactive solute-water aggregates, define
chemically meaningful probe points, and then report time series such as
\[
  V(t),\quad
  |\mathbf{E}(t)|,\quad
  \mathrm{Tr}\,G(t),\quad
  \|G(t)\|_F,\quad
  \eta_G(t).
\]
For microdroplet chemistry, useful probes include bond midpoints,
projected points along putative reaction coordinates, ion-solvation shell
centers, density-isosurface points near adsorbed reactants, and
surface-normal scan points that connect interior and interfacial regions.
These observables can be correlated with bond distances, hydrogen-bond
counts, proton-transfer coordinates, radical-pair separations,
photochemical coordinates, solvent orientation, or vibrational frequency
shifts.  Because ESP, EFIELD, and EFG are evaluated from the same density
and probe coordinates, they also provide a hierarchy of scalar, vector,
and tensor descriptors for testing field-driven mechanistic hypotheses in
microdroplet chemistry and for constructing machine-learning or
charge-model representations of reactive interfaces.

\section{Limitations and Outlook}

The present workflow should be described carefully.  \quick{} can evaluate
ESP, EFIELD, and EFG after each converged energy or gradient call and
exposes that capability through a no-I/O API.  This is sufficient for
trajectory post-processing and for on-the-fly property emission by an
external driver.  It is not yet a complete native AIMD package with
thermostats, restartable dynamics, and automated OEPROP output at every
step.

The next development steps are:
\begin{enumerate}
  \item add production GPU EFIELD and EFG kernels using the existing GPU
    ESP implementation as the template;
  \item benchmark OEPROP overhead on realistic AIMD probe sets, for
    example reactive water clusters, nanodroplet models, and interfacial
    slabs;
  \item add native BOMD hooks that call \keyword{getQuickOEPROP} after
    each accepted step;
  \item extend validation from small fixed water trajectories to
    finite-temperature AIMD trajectories of surface and interior water
    environments;
  \item add small reactive microdroplet benchmarks, such as peroxide,
    thiol, halide, or proton-transfer motifs at finite water interfaces;
  \item document sign, origin, and unit conventions in every trajectory
    output file.
\end{enumerate}
These steps would turn the present validated property and API machinery
into a publishable AIMD electrostatics workflow for local fields and field
gradients in chemically heterogeneous environments, with microdroplet
chemistry as the first application domain.

\section{Data and Software Availability}

The implementation is under development in \quick{}.  Validation inputs,
API drivers, and trajectory-analysis scripts are included in the OEPROP API
and GPU OEPROP validation directories of the development tree.  TODO: add
repository URL, commit hash, archived release DOI, and exact build
instructions for the submitted version.

\begin{acknowledgement}
TODO: Add funding, computing, and contributor acknowledgments.
\end{acknowledgement}

\bibliography{quick_oeprop_aimd_application_note_refs}

\end{document}
